\thispagestyle{empty}\begin{center}\vspace*{-.65in}{\Large\bfseries Carl Friedrich Gauss\par}\bigskip{\Large Werke, Band IV\par}\bigskip{\large Theorie der Kombination von Beobachtungen, die den kleinsten Fehlern ausgesetzt ist\par}\medskip{\large Erster Teil, Zweiter Teil und Ergänzung\par}\end{center}\smallskip\begingroup\small\setlength{\baselineskip}{9.9pt}\noindent Aus dem historischen Latein ins Deutsche KI-übersetzt. Die deutschen Bandtitel, die Anzeigen vom 26. Februar 1821, 24. Februar 1823 und 25. September 1826 sowie die Bestimmung der Genauigkeit der Beobachtungen von 1816 sind in modernem Deutsch wiedergegeben. Dieser Reader enthält die historischen Bandtitel, den förmlichen Titel des Ersten Teils und den vollständigen Ersten Teil, Artikel 1--22, mit beiden Anmerkungen des Autors sowie den förmlichen und den wiederholten Titel des vollständigen Zweiten Teils, Artikel 23--40, mit der Anmerkung des Autors. Es folgt die vollständige Ergänzung mit ihrem förmlichen und ihrem wiederholten Titel, den Artikeln 1--25 und allen vier Anmerkungen des Autors. Auch die vollständigen Anzeigen vom 26. Februar 1821, 24. Februar 1823 und 25. September 1826 aus den Göttingischen gelehrten Anzeigen sind enthalten, gefolgt von der vollständigen Bestimmung der Genauigkeit der Beobachtungen von 1816, Artikel 1--8, mit ihrer handschriftlichen Anmerkung und den Zahlentafeln. Es folgt der vollständige nachgelassene Abschnitt I, \textit{Allgemeine Übersicht der Methode}, zur Bilanz für Witwenkassen, unterzeichnet am 9. Januar 1845 und einschließlich der Anmerkung des Autors, in modernem Deutsch. Es folgt der vollständige Abschnitt II, \textit{Untersuchung des gegenwärtigen Zustands der Professorenwitwenkasse zu Göttingen}, mit dem Vorwort, beiden Abteilungen, allen Anmerkungen des Autors und der am 2. November 1845 unterzeichneten abschließenden Bilanz, in modernem Deutsch. Es folgt der vollständige Abschnitt III, der Kommissionsbericht mit allen vier Fußnoten des Autors und der zugehörigen gesonderten Anmerkung, in modernem Deutsch. Es folgt der vollständige, am 19. Oktober 1851 unterzeichnete Bilanzbericht mit allen neun Positionen und den abschließenden Empfehlungen in modernem Deutsch. Es folgen die zugehörige redaktionelle Einleitung und alle 134 Einträge der Mitgliedertafel. Es folgt die vollständige Leibrentenarbeit mit ihrem Titel, allen zehn Tafeln, Erläuterungen, Rechenbeispielen und Scherings Schlussbemerkungen. Es folgt die vollständige Abhandlung über Abbildungen mit beiden Titeln, dem Vorwort, den Artikeln 1--14, der Fußnote des Autors und den Randbemerkungen in modernem Deutsch. Es folgt die vollständige Abhandlung über gekrümmte Flächen von 1827 mit beiden Titeln, allen 29 Artikeln und der Anmerkung des Autors. Es folgt die vollständige erste Abhandlung über Themen der höheren Geodäsie mit beiden Titeln, dem Vorwort, allen 17 Artikeln, beiden Anmerkungen des Autors und den beigefügten Tafeln in modernem Deutsch. Es folgt die vollständige zweite Abhandlung über Themen der höheren Geodäsie mit beiden Titeln, der Einleitung, den Artikeln 18--35 und der beigefügten Tafel in modernem Deutsch. Vollständig enthalten sind die Anzeigen zu gekrümmten Flächen (1827), höherer Geodäsie (1830, 1843, 1846), Mollweide über Archimedes (1808), Monge und Hachette (1813), Herschel, Mollweide und Kries (1814), Parallelen (1816, 1822) und zum französischen Militärjournal (1830). Es folgen die Ellipsenarbeit und die Carnot-Zusätze (1810) sowie die geometrischen Lösungen von 1823 und 1850 und das vollständige allgemeine Koordinatenverzeichnis mit Scherings Bemerkungen in modernem Deutsch. Es folgen die Abrisse (S.~449--480) mit den dänischen Tafeln in modernem Deutsch. Berichte, Entwürfe und Promemorien zur Gradmessung, die Druckfehlerliste, das Inhaltsverzeichnis und der Druckvermerk beschließen den Band. Das Inhaltsverzeichnis bewahrt die ursprünglichen Seitenangaben.\par\medskip\noindent Historischer Autor: Carl Friedrich Gauss. Der Band wurde 1873 von der Königlichen Gesellschaft der Wissenschaften zu Göttingen herausgegeben. Der Erste Teil wurde am 15. Februar 1821 vorgelegt und 1823 veröffentlicht. Historischer Quelltext und Porträt sind gemeinfrei; Übersetzung und Satz sind KI-produziert. Das historische Original wird gesondert bewahrt.\par\endgroup\clearpage
